\section{Robustness Design}
\label{sec:robustness}

The primary analysis combines the 63-day maximum-drawdown target, the validation-selected model specification, and the portfolio evaluation design described above. The robustness analysis varies one design choice at a time relative to this primary specification, allowing changes in the results to be attributed to a specific alternative rather than to several simultaneous modifications.

\subsection{Predictive Robustness}

Predictive robustness is assessed across alternative definitions and horizons of downside risk. The selected model class or architecture, feature set, and validation-chosen hyperparameters are held fixed, so no new model selection or tuning is performed. The model is re-estimated on the corresponding training and validation observations for each alternative outcome and then evaluated on the test sample. Two alternative 63-day outcomes are considered: downside volatility and the positive magnitude of the worst five-day return. Maximum drawdown is also reconstructed over shorter and longer horizons of 21 and 126 trading days.

All alternative outcomes use the same point-in-time information convention and usability principles as the primary target. The calendar definitions of the training, validation, and test periods remain fixed, but an observation is excluded from a given robustness sample if its alternative target window crosses the next sample boundary. This preserves the chronological separation for each horizon.

\subsection{Investment Robustness}

The portfolio evaluation is examined along six dimensions:

\begin{itemize}
    \item the screen excludes either the most fragile decile or the most fragile quintile;
    \item portfolio weighting changes from equal to market-value weights;
    \item transaction costs of 10, 25, and 50 basis points per unit of one-way turnover are deducted from gross returns;
    \item Newey--West inference for screened-minus-unscreened alpha is repeated with 5, 21, and 63 trading-day lags;
    \item the long--short portfolio is broadened from the extreme deciles (DFRAG-10) to the two least-fragile and two most-fragile deciles (DFRAG-20); and
    \item DFRAG-10 exposure is scaled using its lagged 63-day realized volatility toward a 15\% annualized target, with exposure capped at one so that the strategy is not leveraged.
\end{itemize}

\noindent
Each alternative is compared with the corresponding primary construction without changing the Fragility Score used to rank securities.
